\subsection*{Exercise 1: Universal triangle}
\begin{solution}
A sheafification is a sheaf $\mathcal F^\#$ with $\eta:\mathcal F\to\mathcal F^\#$ such that for every sheaf $\mathcal G$ and every $u:\mathcal F\to\mathcal G$ there is a unique $\bar u:\mathcal F^\#\to\mathcal G$ with $u=\bar u\eta$.  Uniqueness is part of the defining universal property.
\end{solution}

\subsection*{Exercise 2: Uniqueness of associated sheaf}
\begin{solution}
Let $(\mathcal A,\eta)$ and $(\mathcal B,\theta)$ be sheafifications.  Universal properties give $\Phi:\mathcal A\to\mathcal B$ with $\Phi\eta=\theta$ and $\Psi:\mathcal B\to\mathcal A$ with $\Psi\theta=\eta$.  Then $(\Psi\Phi)\eta=\eta$, so uniqueness forces $\Psi\Phi=\id$.  Similarly $\Phi\Psi=\id$.
\end{solution}

\subsection*{Exercise 3: Basis for the espace etale}
\begin{solution}
Every $a\in\mathcal F_x$ is $s_x$ for some $s\in\mathcal F(U)$, so the proposed sets cover $E_{\mathcal F}$.  If $a=s_x=t_x$ lies in two basic sets, equality of germs yields $W\subseteq U\cap V$ with $s|_W=t|_W$.  Then $a\in\widetilde{s|_W}(W)$, which lies in the intersection.  Hence the basis criterion holds.
\end{solution}

\subsection*{Exercise 4: Local homeomorphism}
\begin{solution}
On $\widetilde s(U)$, $\pi$ is bijective with inverse $x\mapsto s_x$.  The inverse sends open basis subsets $W\subseteq U$ to $\widetilde{s|_W}(W)$, while $\pi$ sends such basis subsets to $W$.  Thus both maps are continuous.
\end{solution}

\subsection*{Exercise 5: Local representability}
\begin{solution}
If $\sigma$ is continuous, choose a basic neighborhood $\widetilde s(V)$ of $\sigma(x)$ and shrink so $\sigma(V)$ lies there.  Because $\pi$ is a homeomorphism on that basic open, $\sigma|_V(y)=s_y$.  The converse follows because a map locally equal to continuous maps $\widetilde s$ is continuous.
\end{solution}

\subsection*{Exercise 6: Associated sections form a sheaf}
\begin{solution}
For a compatible family $\sigma_i$ on a cover, define $\sigma(x)=\sigma_i(x)$ when $x\in U_i$.  Compatibility makes this well defined.  It is continuous because continuity is local and satisfies $\pi\sigma=\id$.  Uniqueness is equality of maps on the cover.
\end{solution}

\subsection*{Exercise 7: The canonical morphism}
\begin{solution}
For $V\subseteq U$ and $x\in V$, $(s|_V)_x=s_x$.  Therefore $\widetilde{s|_V}=\widetilde s|_V$, exactly the naturality needed for the maps $\eta_U$ to form a presheaf morphism.
\end{solution}

\subsection*{Exercise 8: Surjectivity on stalks}
\begin{solution}
Let $a\in\mathcal F_x$ and represent it by $s\in\mathcal F(U)$.  Then $\widetilde s\in\mathcal F^\#(U)$ and the germ $(\widetilde s)_x$ maps to $a$.
\end{solution}

\subsection*{Exercise 9: Injectivity on stalks}
\begin{solution}
If two sheafified germs map to the same $a\in\mathcal F_x$, choose local representatives $\widetilde s$ and $\widetilde t$.  Then $s_x=t_x$, so after shrinking $s=t$.  Hence the sheafified germs are equal.
\end{solution}

\subsection*{Exercise 10: Factorization existence}
\begin{solution}
Write $\sigma|_{U_i}=\widetilde{s_i}$.  On overlaps these have equal germs, so $u(s_i)$ and $u(s_j)$ have equal germs.  Since $\mathcal G$ is a sheaf, they agree on overlaps and glue uniquely.  The glued section defines $\bar u_U(\sigma)$; independence follows from common refinements.
\end{solution}

\subsection*{Exercise 11: Factorization uniqueness}
\begin{solution}
If $v,w:\mathcal F^\#\to\mathcal G$ both factor $u$, then on a neighborhood where $\sigma=\widetilde s$ one has $v(\sigma)=u(s)=w(\sigma)$.  These neighborhoods cover, so locality in $\mathcal G$ gives $v=w$.
\end{solution}

\subsection*{Exercise 12: Functoriality}
\begin{solution}
Apply the universal property to $\eta_\mathcal G f:\mathcal F\to\mathcal G^\#$ to obtain $f^\#$.  Both $(gf)^\#$ and $g^\#f^\#$ have composite with $\eta_\mathcal F$ equal to $\eta_\mathcal H gf$, so uniqueness gives equality.  The identity case is identical.
\end{solution}

\subsection*{Exercise 13: Naturality of the unit}
\begin{solution}
The defining equation for $f^\#$ is precisely $f^\#\eta_\mathcal F=\eta_\mathcal Gf$, so the unit is a natural transformation $\id_{\PreSh}\Rightarrow i a$.
\end{solution}

\subsection*{Exercise 14: Stalk map of a sheafified morphism}
\begin{solution}
Taking stalks gives $(f^\#)_x(\eta_\mathcal F)_x=(\eta_\mathcal G)_x f_x$.  Both unit maps are isomorphisms.  Under the identifications they become identities, leaving $(f^\#)_x=f_x$.
\end{solution}

\subsection*{Exercise 15: Sheafifying a sheaf}
\begin{solution}
Since $\mathcal G$ is a sheaf, $\id_\mathcal G$ factors as $r\eta_\mathcal G$.  Then $r\eta=\id$.  The two maps $\eta r$ and $\id_{\mathcal G^\#}$ have the same composite with $\eta$, so uniqueness gives $\eta r=\id$.
\end{solution}

\subsection*{Exercise 16: Idempotence}
\begin{solution}
$\mathcal F^\#$ is a sheaf, so the previous result applies to it and gives $\mathcal F^\#\xrightarrow{\sim}(\mathcal F^\#)^\#$.
\end{solution}

\subsection*{Exercise 17: Constant presheaf on two points}
\begin{solution}
The stalks at $p$ and $q$ are both $A$.  A sheaf on the discrete two-point space has global sections the product of stalks, so $\mathcal P^\#(X)=A\times A$.  A constant global section $a$ maps to $(a,a)$, the diagonal.
\end{solution}

\subsection*{Exercise 18: Bounded continuous functions}
\begin{solution}
A bounded continuous section is continuous, giving a map $\mathcal B^\#\to C^0_X$.  Conversely every continuous $f$ is locally bounded at each point, hence is locally represented by sections of $\mathcal B$ and defines a sheafified section.  These maps are inverse locally and therefore globally.
\end{solution}

\subsection*{Exercise 19: A unit that is not injective}
\begin{solution}
Every global $a\in A$ restricts to zero on a proper neighborhood of each point, so all germs are zero.  Therefore every stalk is zero.  Sheafification has the same zero stalks; a sheaf with every stalk zero is the zero sheaf.  Thus $\mathcal F^\#=0$.
\end{solution}

\subsection*{Exercise 20: Restriction to an open}
\begin{solution}
For $x\in U$, the stalk of $\mathcal F|_U$ is canonically $\mathcal F_x$.  The espace etale of the restricted presheaf is therefore $\pi^{-1}(U)\subseteq E_\mathcal F$, and its continuous local sections are exactly sections of $\mathcal F^\#|_U$.
\end{solution}

\subsection*{Exercise 21: Cofinal basis neighborhoods}
\begin{solution}
Given any neighborhood $U$ of $x$, choose a basis element $B$ with $x\in B\subseteq U$.  Hence basis neighborhoods form a cofinal subsystem of the directed neighborhood category, and filtered colimits over cofinal subsystems agree.
\end{solution}

\subsection*{Exercise 22: Separated presheaf}
\begin{solution}
A presheaf is separated if for every cover, two sections with equal restrictions on all members are equal.  Every sheaf satisfies this by its locality axiom.
\end{solution}

\subsection*{Exercise 23: First plus is separated}
\begin{solution}
If two elements of $\mathcal F^+(U)$ agree locally on a cover, choose representatives by compatible families and pass to a common refinement.  On each refined piece their representatives agree; combining these refinements shows the two cover-data classes are equivalent globally.  Hence locality holds.
\end{solution}

\subsection*{Exercise 24: Plus of a separated presheaf}
\begin{solution}
Let compatible plus-sections be given on a cover.  Represent each by compatible original sections on a refinement.  Refine all overlaps simultaneously.  Compatibility of the plus-sections and separatedness of $\mathcal F$ force the original representatives to agree on common refinements, so their union defines one plus-section on $U$.  Uniqueness follows from separatedness.
\end{solution}

\subsection*{Exercise 25: Two pluses}
\begin{solution}
For arbitrary $\mathcal F$, the first plus $\mathcal F^+$ is separated.  Applying plus again therefore produces a sheaf $\mathcal F^{++}$.  The canonical map $\mathcal F\to\mathcal F^{++}$ has the same factorization property into every sheaf, because compatible local representatives can be glued there uniquely.  Hence $\mathcal F^{++}\cong\mathcal F^\#$.
\end{solution}

\subsection*{Exercise 26: Adjunction}
\begin{solution}
Define $\Phi(h)=h\eta_\mathcal F$.  Given $u:\mathcal F\to\mathcal G$, the universal property supplies the unique $\bar u$ with $u=\bar u\eta$, so $u\mapsto\bar u$ is inverse to $\Phi$.  Compatibility with morphisms in both variables follows from uniqueness of the same factorizations.
\end{solution}

